\section{Stage IV: Multi-Modal and Image-Guided Enhancement}
\label{sec:stage4}

Spatial transcriptomics experiments routinely produce paired histology images (haematoxylin and eosin stains or immunofluorescence) that carry morphological information at sub-micrometre resolution. Stage~IV methods use these complementary modalities, together with spatial coordinate structure, to refine or improve the outputs of earlier stages. Images provide cell-level information (counts, boundaries, morphology) that expression data alone cannot resolve; expression data provide molecular identity that images cannot.

\subsection{Image segmentation as a cell-count prior}

Accurate estimation of the number of cells per spot is a prerequisite for Stages~II and III, yet it is rarely measured directly. Image segmentation fills this gap. Cellpose \citep{stringer2021cellpose} is a generalist deep-learning segmentation model trained on diverse cell morphologies. Applied to H\&E or DAPI-stained images co-registered with spatial transcriptomics slides, it produces per-spot nuclei counts that serve as priors for deconvolution and mapping. Its gradient-flow representation generalises well across tissue types without retraining. StarDist \citep{schmidt2018stardist} uses star-convex polygons to represent nuclear boundaries, offering fast and accurate segmentation particularly suited to densely packed tissues such as tumour cores where touching nuclei are common. Mesmer (DeepCell) \citep{greenwald2022mesmer} is trained specifically on tissue microscopy and provides whole-cell (not just nuclear) segmentation when membrane markers are available, enabling more accurate cell-area estimation.

\subsection{Image-guided deconvolution}

Beyond cell counting, histology images can inform the deconvolution process itself. SegDecon integrates segmentation masks directly into the deconvolution objective: it partitions each spot into sub-regions corresponding to individual segmented cells, then estimates per-cell-type contributions within each sub-region rather than at the whole-spot level. This effectively increases the spatial resolution of deconvolution from spot-level (approximately 55~\textmu m for Visium) to cell-level (approximately 10~\textmu m), bridging the gap between sequencing-based and imaging-based platforms. BayesSpace \citep{zhao2022bulk2space} takes a different approach: it applies a Bayesian spatial prior over a grid of sub-spot positions, using the spatial autocorrelation of expression to enhance resolution by a factor of 3--6 times without requiring image input. When histology is available, BayesSpace can incorporate image features as additional covariates, further improving sub-spot estimation.

\subsection{Integration with imaging-based ST}

A growing experimental design pairs sequencing-based ST (e.g., Visium) with imaging-based ST (e.g., Xenium or MERFISH) on serial sections from the same tissue block. This creates an opportunity for computational integration. Cross-platform anchoring uses the shared gene panel (typically 100--500 genes measured by both platforms) to align coordinate systems and transfer cell-type labels from the imaging platform (which has true single-cell resolution) to the sequencing platform (which has transcriptome-wide coverage). Gene imputation predicts the expression of genes measured only by sequencing-based ST at the single-cell positions identified by imaging-based ST, effectively combining the resolution of one with the breadth of the other. Validation uses imaging-based measurements as ground truth to evaluate and calibrate deconvolution or reconstruction outputs from the sequencing-based data.

This paired-platform strategy is increasingly adopted in atlas-scale projects such as the Human Cell Atlas and offers a practical path toward single-cell-resolved, transcriptome-wide spatial data without requiring computational reconstruction to bear the full burden of resolution enhancement.
